\section{Step-by-step adaptive stitching}
\label{sec:stitch}

\subsection{The exact problem}

Consider an active edge shared by four cells.  In a non-adaptive mesh, the
four cells all have the same size and the edge contributes one segment.  In
an adaptive mesh, the four cells can be a mix of coarse and fine cells.
The edge can cross a boundary between two cells at any coordinate.  The
correct mesh splits the edge at every such boundary and emits one segment
per sub-interval.  The two chunks on either side of a chunk boundary must
produce the same split list; otherwise their rings disagree and a crack
opens.

\subsection{Why a fixed stitching table is not used}

A transvoxel-style table enumerates the finite set of cell configurations
and gives a triangulation for each.  That works when the LoD difference
across a boundary is at most one level and the cell sizes are known.  Here
the difference can be arbitrary (a coarse root cell next to a fine leaf),
the tree is sparse, and the cell sizes are continuous.  A table would have
to enumerate a combinatorial set of nested cells.  The break collector
instead computes the split from the actual cell boundaries, which is
configuration-independent.

\subsection{The recursion, formally}

Let $E$ be an edge span and $[a,b]\subseteq E$ an interval.  Define
\begin{equation}
  B(a,b) = \{\, c_x \mid \text{quadrant cell } c \text{ at } (a+b)/2
  \text{ contributes } c_{\min,x} \text{ or } c_{\max,x} \,\}
\end{equation}
restricted to values strictly inside $(a,b)$ and not within $\epsilon$ of
an existing break.  Then:
\begin{enumerate}
\item append every interior value of $B(a,b)$ to the global break list;
\item for each consecutive pair $a=t_0<t_1<\cdots<t_k=b$, recurse on
  $[t_{i-1},t_i]$ if it is longer than $2\epsilon$.
\end{enumerate}
The recursion terminates because each level inserts at least one break, or
the interval is shorter than $2\epsilon$.  In practice it terminates after
one or two levels: the cell boundaries are powers of two and the interval
is the current edge segment.

\subsection{The GPU stack}

The CPU implementation uses recursion.  The shader replaces it with a LIFO
stack of intervals.  The CPU's exact order is:
\begin{itemize}
\item append all interior breaks first;
\item then recurse into the sub-intervals from left to right.
\end{itemize}
The shader reproduces this by pushing sub-intervals in reverse order:
\begin{lstlisting}[style=glsl,caption={Reverse push preserves the CPU order.},label={lst:push}]
for (int i = localCount - 1; i >= 1; --i)
    pushTask(local[i-1], local[i]);
\end{lstlisting}
Because the stack pops the last pushed interval first, it visits
$[t_0,t_1]$ before $[t_1,t_2]$, and fully expands each before the next.
The global break list is identical to the CPU's.

\subsection{The sample offset}

The quadrant samples are nudged by $\delta$ in the two fixed axes.  The
nudge is not a random jitter: it is a deterministic one-ulp move toward the
quadrant.  The two axes are chosen by \code{edgeAxes}: for an $x$-edge the
fixed axes are $y,z$; for a $y$-edge they are $x,z$; for a $z$-edge they
are $x,y$.  The four quadrant signs are fixed by the \code{quadrantSigns}
table:
\begin{center}
\begin{tabular}{@{}c|cccc@{}}
\toprule
axis & $q_0$ & $q_1$ & $q_2$ & $q_3$\\
\midrule
$x$ & $(-1,-1)$ & $(-1,+1)$ & $(+1,+1)$ & $(+1,-1)$\\
$y$ & $(-1,-1)$ & $(+1,-1)$ & $(+1,+1)$ & $(-1,+1)$\\
$z$ & $(-1,-1)$ & $(-1,+1)$ & $(+1,+1)$ & $(+1,-1)$\\
\bottomrule
\end{tabular}
\end{center}
The actual table is flattened as
\code{QS[axis*8 + quadrant*2 + {0,1}]}; its values are shown above.

\subsection{Break capacities and overflow}

The fixed capacities are a hard contract.  A scan cell with more than 64
breakpoints, or an interval with more than 32 local breaks, or a recursion
stack deeper than 256 intervals is not expected for the engine's cell
sizes, but it is still possible for a pathological SDF.  The shader does
not truncate:
\begin{itemize}
\item \code{noteOverflow()} sets \code{gOverflow};
\item the cell stops emitting;
\item \code{overflowCount} is incremented;
\item the host sees the counter and fails the chunk.
\end{itemize}
The host retries the chunk after the tree edit or reports it as skipped.
It never publishes a partial index stream.  This is an important
difference from a ``best effort'' GPU mesher: a silent truncation would be
a hole in the terrain that is hard to diagnose.

\subsection{The global break list}

The global list is sorted and deduplicated by insertion sort and a linear
unique pass.  The lists are small (at most 64 entries), so an $\Oh(n^2)$
sort is faster than a library sort and, more importantly, has a
deterministic order.  Equality uses the edge epsilon:
\begin{equation}
  |b_i-b_j|\le \epsilon \quad\Rightarrow\quad \text{duplicate}.
\end{equation}
The sort is stable for equal keys because the insertion sort only moves
strictly greater values.  This matches the CPU's \code{std::sort} behaviour
for these list sizes and keeps the emitted triangle order deterministic.

\subsection{Owner and quadrant sampling}

For each sub-segment the shader samples the four quadrants at its
midpoint.  It does not reuse the quadrant cells from the break collection:
the midpoint of a sub-segment is a different position and may fall in a
different cell.  The owner is selected by the ordering rule of
\Cref{sec:cpuwalk}.  The owner's interpolated field at the segment
endpoints decides the winding.  If the owner's interpolation degenerates
(both endpoints get the same sign because a zero lies exactly on a corner),
the producing cell's endpoint signs are used instead.  Those signs are
authoritative because the producing cell only scans edges with a sign
change.

\subsection{Consequences}

The stitching step is what makes it legal to publish every LoD rung and let
the renderer choose one per region.  Without it, the mesh would be
crack-free only within a single rung; at rung transitions the surface would
have T-junctions and visible seams.  With it, the coarse and fine sides of
a boundary agree on the edge split and on the owner, so their duplicated
boundary quads have identical vertex order and identical triplanar plane.
This is also why the chunk-border guard in the simplifier is necessary:
the stitching can resolve interior LoD changes, but it cannot resolve a
difference in the chunk's own boundary cells because the two chunks never
share a node table.
